\section{Least bounded search and finite choice}
\label{sec:search}

Bounded search is the canonical way to turn a decidable finite existence statement
into a primitive-recursive witness.  It is also the basic deterministic selector used
inside demand passes, viability tables, finite nets, and certificate checkers.  The
repository already contains a checked object implementation in
\code{IPRADerivations.Core.Search}; we state the abstract theorem because it is a
building block of the normal form.

\subsection{Definition by primitive recursion}

Let \(p(i,x)\in\B\) be a primitive-recursive Boolean predicate, with the convention
that \(p(i,x)=1\) means success.  Define
\[
 \mu_p(0,x)=0
\]
and, for \(N\ge0\),
\[
 \mu_p(N+1,x)=
 \begin{cases}
   \mu_p(N,x),&\mu_p(N,x)<N,\\
   N,&\mu_p(N,x)=N\ \text{and}\ p(N,x)=1,\\
   N+1,&\mu_p(N,x)=N\ \text{and}\ p(N,x)=0.
 \end{cases}
\]
Every branch is expressible by equality tests and selectors, so \(\mu_p\) is primitive
recursive.  Equivalently, using a sign convention in which zero means success, one may
use the recurrence already implemented in the repository.

\begin{theorem}[Least bounded search]
\label{thm:bounded-search}
The function \(\mu_p\) satisfies, and \(\iPRA\) proves:
\begin{enumerate}[label=(\roman*),leftmargin=2.4em]
\item \(\mu_p(N,x)\le N\);
\item if \(\mu_p(N,x)<N\), then \(p(\mu_p(N,x),x)=1\);
\item if \(i<\mu_p(N,x)\), then \(p(i,x)=0\);
\item if \(i<N\) and \(p(i,x)=1\), then
      \(\mu_p(N,x)<N\) and \(\mu_p(N,x)\le i\).
\end{enumerate}
\end{theorem}

\begin{proof}
Use one induction on \(N\) with a joint quantifier-free invariant encoding (i)--(iii)
for an arbitrary comparison index \(i\).  The base case is immediate.  For the step,
there are two cases.  If the old search value is below \(N\), the recurrence freezes
it; all three clauses follow from the induction hypothesis.  If the old value equals
\(N\), inspect the Boolean \(p(N,x)\).  A success freezes the new value at \(N\); a
failure moves it to \(N+1\).  In either case bounded minimality follows by combining
the old minimality with the value at \(N\).  Clause (iv) is then derived by
contradiction: if \(i<\mu_p(N,x)\), clause (iii) contradicts \(p(i,x)=1\); otherwise
\(\mu_p(N,x)\le i<N\), and clause (ii) supplies success.  All case distinctions are on
numerical equalities or Boolean values, so the invariant is quantifier-free.
\end{proof}

\subsection{Finite choice is harmless when its bound is explicit}

A common constructive concern is the use of ``choose an \(i\)'' from a finite set.
In this setting the concern is not finite choice itself.  Once a bound \(N\) and a
Boolean verifier are available, \cref{thm:bounded-search} gives a canonical least
witness.  The real questions are:

\begin{itemize}
\item where the bound \(N\) came from;
\item whether the verifier captures the intended mathematical property; and
\item whether choosing the least encoded witness is extensionally legitimate.
\end{itemize}

The first question leads to producer boundaries and the cardinality-to-cutoff barrier.
The second is an object-proof obligation.  The third motivates relational outputs:
least choice is computationally available but may not respect equality of represented
objects.  One should therefore choose least only after proving that the downstream
observable is independent of the enumeration, or use least choice merely as an
internal path selector whose output is later erased by a stability theorem.

\subsection{Bounded searches over coded tuples}

Pairing and simultaneous recursion allow a search over any fixed finite product to be
reduced to one numerical interval.  For example, to search
\((i,j)\) with \(i<I\) and \(j<J\), use the code \(k=iJ+j<IJ\), recover
\(i=k\mathbin{/}J\) and \(j=k\bmod J\), and Booleanize the range conditions.  More
complex bounded records may be encoded by nested pairing.  The size of the search
space may be enormous; \(\iPRA\) is concerned with primitive recursiveness and proof,
not practical efficiency.  Nevertheless, the normal form encourages structured
loops and dynamic programming when they yield smaller invariants.

\subsection{The selector/finality distinction}

Suppose a finite cloud \(C_d(x)\) contains several valid approximations at precision
\(d\).  Bounded search yields the least valid code \(c_d\).  This proves existence of a
primitive-recursive selector for each finite level.  It does \emph{not} prove that
\((c_d)_d\) is coherent, convergent, or presentation invariant.  Those are separate
statements about compatibility edges, diameter bounds, and motion budgets.  We will
therefore use bounded search in two roles:

\begin{enumerate}[label=(\alph*),leftmargin=2.4em]
\item as an exact selector when equality and ordering of candidates have mathematical
meaning; and
\item as a deterministic implementation detail inside backward viability, followed by
a proof that the final observable does not depend on that detail.
\end{enumerate}

This distinction prevents a finite-choice theorem from being silently upgraded to a
choice principle for infinite represented objects.
